\section{Certified errors in the discrimination controls}
\label{sec:controlstability85}
The preceding lower bounds first concern ideal known-pair controls.
Their stability has a direct finite formulation.  An error bound on a
control channel is a bound in unhalved diamond norm, valid on every
input including its memory reference.  It is not a bound measured only
on one calibration state.

\begin{theorem}[A finite control-error ledger]\label{thm:controlstability85}
Use the code and recovery from Lemma~\ref{lem:curvelogical85}, with
gap $\Delta$ and remainder constant $\kappa$.
Suppose an implemented initial logical state differs from the ideal
one in trace norm by at most $e_0$.  At cycle $r$, suppose its
implemented logical channel differs from the corresponding ideal
corrected channel by at most $\nu_r$ under each hypothesis.
Suppose the implemented final readout differs from the ideal
Helstrom readout by at most $e_f$ in diamond norm.
For every integer $m\le N$ and $|t|\le a_0$, the actual classical
output laws obey
\begin{equation}\label{eq:controlledger85}
 \|\widetilde p_{0,m}-\widetilde p_{t,m}\|_1
 \ge\left(2|\sin(mt\Delta/2)|-m\kappa t^2
              -2e_0-2e_f-2\sum_{r=1}^m\nu_r\right)_+.
\end{equation}
Here one may choose the Helstrom readout for the ideal unrotated and
logical-unitary states; it depends on the known pair and on $m$.
The cycle error may be bounded by the sum of the encoding and recovery
errors, because the unknown measurement itself is a channel.

In particular, put $z=\min\{1,N|t|\Delta\}$ and choose $m,t_0$
as in the proof of Theorem~\ref{thm:curveclassification85}.
For $0<|t|\le t_0$, the sufficient conditions
\begin{equation}\label{eq:controlprecision85}
 e_0+e_f\le z/(16\pi),\qquad
 \nu_r\le |t|\Delta/(8\pi)\quad(1\le r\le m)
\end{equation}
imply
$\|\widetilde p_{0,m}-\widetilde p_{t,m}\|_1\ge z/(2\pi)$.
All errors are deterministic uniform channel bounds; no independence
of control errors is required.
\end{theorem}
\begin{proof}
Choose the ideal two-outcome Helstrom readout for the equal logical
superposition and its image under $e^{-imt\Gamma_L}$.
Its classical $\ell^1$ separation is $2|\sin(mt\Delta/2)|$.
Replacing the logical unitary by the actual ideal corrected channel
loses at most $m\kappa t^2$, by channel telescoping and
Lemma~\ref{lem:curvelogical85}.  Under either hypothesis,
telescoping the implemented preparation, cycles and readout changes
its output law in $\ell^1$ by at most
$e_0+e_f+\sum_r\nu_r$.  The triangle inequality proves
\eqref{eq:controlledger85}.  Replacing encoding and recovery around
one unchanged unknown channel costs no more than the sum of their
diamond errors, which proves the stated decomposition of $\nu_r$.

For the indicated choice set $x=m|t|\Delta$, so $z/2\le x\le1$.
The ideal sine term minus curvature is at least $3x/(2\pi)$.
Conditions \eqref{eq:controlprecision85} give
$e_0+e_f\le x/(8\pi)$ and $\sum_r\nu_r\le x/(8\pi)$.
The total loss is therefore at most $x/(2\pi)$, leaving
$x/\pi\ge z/(2\pi)$.
\end{proof}

For a unitary orbit, Lemma~\ref{lem:correctedchannel84} gives the
same result with $\kappa=4\|G\|_{\mathrm{op}}^2$ and the original
support-code gap.  An error allowance fixed independently of the
angle need not meet \eqref{eq:controlprecision85} as $t\to0$.
The theorem does not certify a particular piece of hardware, synthesize
a gate circuit, or allow target-dependent advice in common learning.
It specifies exactly which certified approximation errors suffice
to preserve the finite-use lower bound.

\begin{proposition}[Exact first-order decisions]\label{prop:curveexact85}
For Gaussian-rational $E$ and a Gaussian-rational Hermitian zero-sum
$H$, two-sided first-order realizability and the branch of
Theorem~\ref{thm:curveclassification85} for $H\ne0$ are decidable
in polynomial represented-input bit complexity.  The decision gives
$\Gamma_E(H)$, its projection on $\mathcal W_E$, and its orthogonal
residual.  A zero tangent is reported as a higher-order question,
not as stationarity of an unspecified curve.
\end{proposition}
\begin{proof}
The exact support projectors in Proposition~\ref{prop:supportexact84}
are rational.  Test $Q_jH_jQ_j=0$ for every $j$ and compute
\eqref{eq:curvegenerator85}.  Necessity and sufficiency of the test
follow from Lemma~\ref{lem:curvegenerator85} and the normalized-factor
construction in Lemma~\ref{lem:curvelogical85}: for any such tangent
it gives a real-analytic measurement curve with that derivative.
Real coordinate rank selection and the Hilbert--Schmidt Gram system
then project $\Gamma_E(H)$ exactly.  Rational matrix arithmetic and
fraction-free elimination have the bit bounds already proved in
Proposition~\ref{prop:supportexact84}.  Only these finite operations
are performed.  A numerical curvature bound for an unspecified
curve, an optimal tester, and a compiled recovery are not outputs
of this decision.
\end{proof}
